\section{系统测试与多维度实测验收}

为了避免“单测全绿但线上失灵”的常见盲区，本项目按照《验收设计》（A01--A37，完整覆盖 R01--R24）建立了五层递进验收体系：\textbf{单元与契约回归 $\rightarrow$ 知识库容量与复杂样本抽查 $\rightarrow$ 真实大模型端到端规划与偏好对照 $\rightarrow$ 全页面浏览器与无障碍验收 $\rightarrow$ Linux 容器与空卷恢复}。所有真实在线验收均保留原始请求、草案 JSON、任务日志与独立复算报告，对未达标或尚待核对的问题如实记录，绝不粉饰为通过。

\subsection{自动化测试与静态检查基线}

截至 2026 年 11 月 22 日最终回归，前后端全量自动化测试与代码质量检查结果汇总于表~\ref{tab:auto-baseline}。

\begin{table}[H]
  \centering
  \caption{自动化测试与构建质量基线}
  \label{tab:auto-baseline}
  \renewcommand{\arraystretch}{1.28}
  \begin{tabularx}{\textwidth}{@{} >{\sffamily\bfseries}p{3.4cm} >{\sffamily}p{3.0cm} Y @{}}
    \toprule
    验证维度 & 实测结果 & 核心覆盖范围 \\
    \midrule
    后端单元测试（\texttt{pytest}） & \textbf{349 项通过}（14.94 s） & 配置与脱敏、SQLite v6 迁移、四家搜索回退与 MCP、高德限流、动态循环与检查点重试、事实引用与跨年季节核验、局部编辑与旧版本冲突、离线备份恢复、真实 HTTP 断连与进程强杀恢复（已隔离代理）。 \\
    Python 静态检查（\texttt{ruff}） & \textbf{全部通过}（130 文件） & 代码规范 \texttt{ruff check} 与格式一致性 \texttt{ruff format --check} 零警告。 \\
    前端类型与单测（\texttt{tsc / Vitest}） & \textbf{60 项通过}（5 文件） & 金额分精度、人数与时间约束、跨午夜边界、费用依据保留与业务工具函数。 \\
    前端生产构建（\texttt{Vite}） & \textbf{构建通过}（244 ms） & 产出静态资源并经容器构建验证 SHA-256 一致。 \\
    \bottomrule
  \end{tabularx}
\end{table}

\subsection{知识库容量压测与真实复杂资料抽查}

知识库分别通过了“规模容量压测”与“多形态真实文档抽查”两轮独立验收：

\begin{itemize}
  \item \textbf{100 份中文资料容量与强杀恢复（A21）}：实际 HTTP 导入 50 份中文文本与 50 份中文 PDF（合计 28.32 MB，含 190 万字符长文、200 页 PDF 与恰好 20 MiB PDF），全部进入 \texttt{ready} 并生成 4,140 个片段，总导入耗时 20.20 秒；100 次专属词 HTTP 检索中位数 \textbf{5.29 ms}（P95 为 7.16 ms），进程组峰值 RSS 459.64 MiB。第 101 份（\texttt{409}）、20 MiB + 1 字节（\texttt{413}）、201 页 PDF 与超 200 万字符文本均被准确拒绝；解析长文时强杀进程组重启后标记 \texttt{failed/import\_interrupted} 且零脏索引。
  \item \textbf{五类真实形态样本与 SSRF 防护（A15 / A16 / A25）}：在独立空库中导入 26 页中文排版 PDF（43 片段，“网格”精准命中第 19 页）、2 页纯图片扫描 PDF（准确进入 \texttt{failed/pdf\_no\_text}，0 片段并提示提供文字版）、3,083 字西安攻略长文（67 片段，“碑林”命中第 5 段）、陕西历史博物馆官网 URL（21 片段）与秦始皇帝陵博物院收藏链接（126 片段）；额外提交本机回环 URL 被准确拦截为 \texttt{failed/private\_url}（0 片段）。原文下载字节/SHA-256 全部一致，删除中文 PDF 后 43 条片段与 FTS5 倒排记录同步归零，原文件目录移除，再次检索无残留。
\end{itemize}

\subsection{西安五日游真实端到端规划、局部编辑与恢复}

本轮正式端到端验收直接读取更新后的 \texttt{.env}（DeepSeek \texttt{deepseek-flash} + 高德 Web 服务 + 四家搜索顺序 \texttt{Exa $\rightarrow$ DeepSeek $\rightarrow$ Tavily $\rightarrow$ AnySearch}，超时 120 秒）。显式服务检查确认模型、高德地图与 Exa 匿名 MCP 均一次连通成功。

\textbf{输入条件}：西安，2026-12-05 至 2026-12-09（5 天），2 名成人，全团预算 3,500 元；首日 10:00 抵达西安北站、末日 18:00 由西安北站离开；住宿自理于南稍门或钟楼附近；喜欢历史文化、博物馆与老街区，愿意提前预约但尚未预约；午晚餐为自由用餐时段；同时显式关联 1 份西安参观参考笔记（内含独有说明“赏碑时想看石刻笔画与拓片工艺”）。

\begin{table}[H]
  \centering
  \caption{西安五日游真实生成日程概要（版本 1 / 版本 2 对照）}
  \label{tab:xian-5days}
  \renewcommand{\arraystretch}{1.26}
  \begin{tabularx}{\textwidth}{@{} >{\sffamily\bfseries}p{2.0cm} Y @{}}
    \toprule
    日期 & 实际排程（版本 1；括号内为局部编辑后版本 2 变化） \\
    \midrule
    12-05（首日） & 10:00--10:30 西安北站抵站；11:30--12:30 自由午餐（钟楼附近）；13:00--14:00 西安钟楼；14:30--16:30 北院门历史文化街区；17:30--19:00 自由晚餐（北院门附近）。 \\
    12-06（次日） & 09:00--12:30 秦始皇帝陵博物院\textbf{（版本 2 缩短 15 分钟至 12:15，门票 \textyen{}240 升为已核实）}；12:30--13:30 自由午餐（景区附近）；18:00--19:30 自由晚餐（钟楼附近）。 \\
    12-07（三日） & 09:00--10:30 西安城墙永宁门；11:00--13:00 西安碑林博物馆；13:00--14:00 自由午餐（碑林附近）；18:00--19:30 自由晚餐（钟楼附近）。 \\
    12-08（四日） & 09:00--11:30 陕西历史博物馆；12:00--13:00 自由午餐（陕历博附近）；14:00--16:00 大雁塔；17:30--19:00 自由晚餐（大雁塔附近）。 \\
    12-09（末日） & 10:00--11:30 易俗历史文化街区；12:00--13:00 自由午餐（易俗街区附近）；15:30--18:00 西安北站安检候车。 \\
    \bottomrule
  \end{tabularx}
\end{table}

\textbf{实测结论与关键指标}：
\begin{enumerate}
  \item \textbf{规划质量与零硬冲突}：版本 1 用时 1,566.87 秒（8 次决策、22 次工具调用），产出 5 天共 19 项活动，19 项全部附有独立生成的简洁介绍；9 项餐饮均为自由用餐且无指定店名；非餐饮活动按 POI 统计零重复游览（仅首末日抵离复用西安北站）。参考笔记中独有的“赏碑时想看石刻笔画与拓片工艺”被成功吸收至规划意图。Phase 1 的规则与提示词修复全面见效：\textbf{首尾日时段覆盖误报、枢纽站营业时间误报以及步行/接驳冲突全部清零}，版本 1 核验状态为 \textbf{\texttt{needs\_review}（0 条 \texttt{conflict}，38 条 \texttt{unknown}）}，56 条事实中 15 条为官方 \texttt{supported}、30 条 \texttt{estimated}、11 条 \texttt{unknown}、0 条冲突。
  \item \textbf{定点局部编辑严格保真}：针对第二天将秦始皇帝陵博物院时长由 210 分钟缩短为 195 分钟，耗时 260.68 秒生成版本 2，状态仍为 \textbf{\texttt{needs\_review}（0 条 \texttt{conflict}，36 条 \texttt{unknown}）}。逐字段比对确认：\textbf{未受影响的其余 4 天共 16 项活动、12 段交通路线与 45 条事实依据 100\% 原值保持不变}，历史版本 1 只读回读完全一致；同时第二天重查事实使兵马俑两张成人票（\textyen{}240）从估算晋升为已核实，全团已知费用由 \textyen{}170.00 提升至 \textyen{}410.00（估算由 \textyen{}1,594.00 同步调整为 \textyen{}1,354.00，小计恒为 \textyen{}1,764.00，7 项未知，复算误差为 0）。
  \item \textbf{停服校验、备份与异目录恢复}：停服后依次执行 \texttt{verify}、\texttt{backup}、向新空目录 \texttt{restore} 并再次 \texttt{verify}，原库与恢复库均确认 Schema v6、1 行程、2 任务、2 版本、1 资料、1 片段、0 孤立文件，通过 HTTP 读回的两个版本逐字段一致。
\end{enumerate}

\subsection{个人偏好真实大模型 A/B/C 对照验证}

为了证明自然语言长期偏好不仅能“存得住”，而且在真实模型推理中“\textbf{真的生效、会被本次明确要求覆盖、绝不偷偷篡改档案且不展示匹配理由}”，我们在独立实例执行了三阶段真实对照实验（西安 3 日，2 人，预算 2,400 元，长期偏好统一设为：“每天不早于 09:30 开始游览；偏爱博物馆与古建，不逛夜景和商业街；单段步行尽量不超过 15 分钟”），核心对比见表~\ref{tab:pref-ab}。

\begin{table}[H]
  \centering
  \caption{个人偏好真实大模型 A/B/C 对照实验结果}
  \label{tab:pref-ab}
  \renewcommand{\arraystretch}{1.28}
  \begin{tabularx}{\textwidth}{@{} >{\sffamily\bfseries}p{2.0cm} >{\sffamily}p{3.1cm} Y @{}}
    \toprule
    实验组 & 额外输入与操作 & 实测可观测结果与规则判定 \\
    \midrule
    A 组（正向生效） & 仅提供抵离与住宿，不写节奏或景点偏好 & 意图完整识别全部 4 项长期习惯；完整日 09:30 起步，选入碑林、陕历博、西安博物院与永宁门，\textbf{零商业街与夜游}；15 项介绍 0 条偏好匹配理由；长期偏好原文未变。另真实暴露 1 处模型将钟楼至兵马俑 41.2 km 误排步行的失误，被规则引擎诚实拦截为 \texttt{blocked}（\texttt{travel\_time} 冲突）。 \\
    B 组（本次优先覆盖） & 保持长期偏好不变，本次明确增加“节奏紧凑，必须安排大唐不夜城夜游” & 意图显式记录“按本次要求执行夜游”，并成功排入 \textbf{12-08 19:10--21:00 大唐不夜城夜游}；行程状态为 \textbf{\texttt{needs\_review}（0 硬冲突）}；20 项介绍 0 条偏好匹配理由；\texttt{GET /api/preferences} 原文与版本号完全未变。 \\
    C 组（确认更新与乐观锁） & 先用旧版本 1 提交更新，再用正确版本 2 主动修改并重启 & 旧版本准确返回 \textbf{HTTP 409（\texttt{revision\_conflict}）} 且读回值未变；正确版本升至版本 3 后服务重启读回与离线 \texttt{travel-data verify} 全部通过。 \\
    \bottomrule
  \end{tabularx}
\end{table}

\subsection{全页面 Playwright 浏览器验收与 Linux 容器复验}

在交互层与交付层，系统同样完成了全量端到端验证：
\begin{itemize}
  \item \textbf{Playwright 28 组浏览器流程与无障碍检查（用时 39.65 秒）}：覆盖模块 1（行程新建边界、非法日期/2 天/8 天拒绝、精确到分/零预算/留空、时间约束与跨午夜预约、表单内参考资料粘贴/上传/链接/取消关联/失败重试）、模块 2（参考资料五类管理、失败重新导入、中文检索定位、原文下载、分页与删除清零）、模块 3（偏好保存、清空与 409 并发冲突重载）、模块 4（四家搜索顺序含 \textbf{DeepSeek} 展示、密钥/匿名/禁用状态区分、缺配置提示与脱敏校验）以及模块 5（设置导航、历史版本只读、局部编辑、补充弹窗）。桌面 1440px 与手机 390px/320px 均实现 \textbf{0 横向溢出}、\textbf{0 页面异常}、\textbf{0 意外控制台错误}；对 38 个页面及弹窗状态执行 axe-core 审计（WCAG 2.0/2.1/2.2 A/AA），\textbf{违规数为 0}。
  \item \textbf{Linux Docker 容器构建与空卷恢复（用时 35.81 秒）}：基于最新源码与前端产物构建 \texttt{linux/amd64} 镜像，启动后健康检查返回 \texttt{schema\_version: 6}，容器内 JS/CSS 构建资源与本机构建的 SHA-256 逐一相等；核验 UID/GID 为 \texttt{10001:10001}、镜像内无 \texttt{.env}/源码/\texttt{pytest}、宿主机端口仅绑定 \texttt{127.0.0.1}；写入数据后强制替换容器数据完整保留，停服校验、打包备份并向新空命名卷恢复后，行程、偏好、原资料与中文全文索引全部一致通过。
\end{itemize}
